\begin{abstract}
A language model can propose a plausible correction for a failed attempt without knowing whether it will help the agent finish. Even when an intervention works, deciding what the policy should learn from it is a separate problem. We introduce \method{}, a research harness that connects failure analysis to policy learning through same-state, paired environment tests. Only interventions that repeatedly improve terminal completion become positive training evidence. A dual-view representation pairs a deployable action with training-only supervision for subgoals, expected progress, and recovery; evaluation then removes external memory and intervention. An ALFWorld study with Qwen2.5-3B-Instruct exposes a gap between these stages. Five development rounds produce one stable beneficial intervention among 21 source-state experiments. A subsequent executable-strategy validation finds four among five selected states, but these successes are completed by the intervention programs themselves. Training on their eight dual-view examples leaves scaffold-free selection success unchanged at 3 of 355 tasks, with one gain and one regression. The result separates useful intervention from acquired capability: environment verification can establish what works at a state without establishing that a small policy has learned to reproduce it. \method{} makes that distinction observable and identifies incomplete action coverage and persistent execution errors as concrete targets for further study.
\end{abstract}
